% Постановка проблеми.
% Аналіз останніх досліджень та публікацій.
% Формулювання цілі статті.
% Виклад основного матеріалу.
% Висновки.
% Список літератури англійською мовою за стандартом APA Style (використання транслітерації при переклад

\begin{center}
  \textbf{Вступ}
\end{center}

Сертифіковані медичні прилади зазвичай вимагають значних фінансових інвестицій. Кожен із ним має відповідний процес
вишколювання персоналу. Проте, значну частину знань можна отримати без реального приладу, а у симуляції. Для цього
будують симулятори органів людського тіла. Використання такого підходу здешевлює процес навчання і полегшує логістику
вишколу. Звісно він не замінить реальні прилади, але суттєво полегшує вишколювання спеціалістів.

Моделювання процесів теплообміну у тканинах людського тіла можна описати кількома моделями
\cite{kutz-zhu-heat-transfer-biological-systems}:
\begin{itemize}
  \item Рівняння Пеннеса \cite{thermal-process-tissue-pennes} -- у модельованих м'яких тканинах береться усереднений
        ефект від потоку крові через об'єм.
  \item Моделювання головних судин, які впливають на температурні процеси органу чи органів, які досліджуються.
\end{itemize}
Із часом з'являються досконаліші інструменти, які суттєво полегшують створення судинних моделей. Проте вони є
складнішими, але дають точніші результати. Модель описана рівнянням Пеннеса є достатньо для багатьох симуляцій і її 
регулярно уточнюють \cite{thermal-process-tissue-pennes-explained}, \cite{elsapaa2023thermoelastic}.

На практиці рівняння Пеннеса використовується для моделювання теплової терапії. Вона використовується як допоміжна разом
із хіміотерапією і променевою терапією при лікуванні ракової пухлини. Теплова терапія нагріває пухлину за допомогою 
радіочастотної або мікрохвильової енергії \cite{elsapaa2023thermoelastic}. 

Лазерна абляція використовується для усунення доброякісних і злоякісних утворень у тканинах людського тіла за допомогою
лазера. У праці \cite{munshid2025laser} описано підхід для моделювання цієї процедури. Інший метод, який можна описати
рівнянням Пеннеса це -- кріотерапія \cite{fahmy2025fractional}. Він використовує наднизькі температури для заморожування
та руйнування пухлин.

Для отримання розв'язку рівняння Пеннеса використовуються аналітичні методи \cite{trefilov2025generalizedbioheat}. Вони
дозволяють знайти точне представлення розв'язку. На противагу ним існують чисельні методи. Найпоширеніші серед них це 
метод скінчених елементів (МСЕ) \cite{bahvalov-et-al}. Ще одним чисельним методом, який застосовується до цих задач є
метод граничних елементів (МГЕ) \cite{brebbia-bem}. Відносно молодим чисельним методом є метод приграничних елементів 
(МПГЕ) \cite{grytsko-cobem}. Він належить до сімейства МГЕ, проте має кілька переваг: 
\begin{itemize}
  \item Інтегрування відбувається по приграничній області.
  \item Метод має додатковий степінь вільності, а саме ширину приграничної зони. Він дозволяє покращувати точність 
  обчислень.
\end{itemize}
Як бачимо, для розв'язування задач на базі рівняння Пеннеса використовуються різноманітні методи. Таким чином, 
застосування більш ефективніших підходів для отримання розв'язків є актуальним напрямом досліджень. 

\textit{Об’єкт дослідження} -- моделювання процесів теплопровідності в органах людського тіла. 

\textit{Предмет дослідження} -- методи для розв'язування задач на базі рівняння Пеннеса. 

\textit{Мета дослідження} -- аналіз методів та засобів для моделювання процесів теплопровідності в органах людського
тіла.

\begin{center}
    \textbf{Постановка проблеми}
\end{center}

Поширені методи для розв'язування задач для рівняння Пеннеса не завжди є ефективними. Можна покращити точність або 
зменшити час виконання методу завдяки перевагам МПГЕ. Процес моделювання поширення тепла в органах людського тіла може
використовуватися у медичному обладнанні або для навчання медичного персоналу в умовах наближених до реальних. Тому 
розробка оптимальних методів для обчислення теплового режиму є нагальною проблемою.
Згідно із вищенаведеними проблемами, основною метою дослідження є застосування МПГЕ для розв'язування задач для рівняння
Пеннеса і порівняння отриманих результатів із відомими аналітичними розв'язками, а також із результатами із МСЕ.   

\begin{center}
  \textbf{Аналіз останніх досліджень та публікацій}
\end{center}

\cite{ostadhossein2023bioheat} розв'язують одновимірні лінійні і нелінійні задачі для рівняння Пеннеса. Ці рівняння
описують вплив різкий перепад температур на поверхні людської шкіри. Коефіцієнти питомої теплоємності, теплопровідності
і перфузії крові розглядаються як лінійні функції від температури. Для розв'язування задачі використовується метод
скінчених різниць (МСР) і метод скінчених об'ємів (МСО). Робляться висновки, що коефіцієнт теплопровідності і перфузія
крові, залежні від температури, мають суттєвий вплив на модель із різким перепадом температур на поверхні шкіри, а
питома теплоємність може залишатися сталою.

Автори статті \cite{gangadhara2024stability} використовують МСЕ для аналізу критеріїв стабільності для рівняння переносу
біологічного тепла у людському тілі. МСР використовується для дискретизації похідної по часу. Робляться висновки про 
критерії стабільності. 

Multigrid Waveform Relaxation Method (WRMG) використовується у поєднанні із МСР у \cite{santiago2023multigrid} для 
розв'язування рівняння Пеннеса. Автори статті проаналізували інші наукові праці і дійшли висновку, що цей метод ще не 
використовувався для розв'язування рівняння теплопровідності у людському тілі. Цей метод оптимізований для паралельного 
виконання. Автори використали цей метод для розв'язування на двовимірній академічній задачі для розподілу температури на
людській шкірі.

Праця \cite{munshid2025laser} порівнює три методи для моделювання теплових процесів які виникають у людській шкірі
внаслідок впливу лазерного випромінювання. Зокрема розглядаються МСЕ, Bacterial Foraging optimization (BFO) і
аналітичний розв'язок. Автори формулюють задачі, для перевірки методів. Робляться висновки і рекомендації на основі
порівняльного аналізу цих методів.

Одним із популярних симуляторів із відкритим кодом є iMSTK \cite{imstk_reference}. У \cite{imstk_ablation_liver} 
описується підхід із використанням цього симулятора для моделювання мікрохвильової абляції у печінці. iMSTK використовує
Position Based Dynamics (PBD) для моделювання динаміки фізичних тіл під дією зовнішніх сил. Також у цьому симуляторі 
можна використовувати МСЕ для моделювання фізичних процесів. Автори статті розв'язують рівняння Пеннеса для моделювання
впливу мікрохвильової абляції на печінку людини. Завдяки засобам iMSTK вони інтегрують цю інформацію з іншими модулями
системи щоб створити реалістичну симуляцію для медичних працівників чи студентів. 

\cite{transmission_line_model} досліджують взаємодію між хвилями і м'якими тканинами людського тіла. Пропонується 
математична модель в якій вираз на основі питомого коефіцієнта поглинання електромагнітної енергії додається як зовнішнє
джерело тепла до рівняння Пеннеса. Це дозволяє змоделювати збільшення температури м'якої тканини під впливом 
електромагнітного випромінювання. Автори використали transmission line model (TLM) для розв'язування задачі і порівняли 
отриманий розв'язок із МСР.

У праці \cite{trefilov2025generalizedbioheat} досліджують модель лазерної абляції. Для цього автори удосконалюють
модель, на основі рівняння Пеннеса для опису впливу зовнішнього випромінювання на м'які тканини людини. Зокрема модель
додає більш реалістичне представлення теплообміну між органами і кровотоком з урахуванням нерівномірного розподілу 
тепла. Отримані аналітичні розв'язки узгоджуються із числовими отриманими в програмному пакеті COMSOL, який використовує 
МСЕ.

Стаття \cite{tasnim2024thermal}...

\begin{center}
  \begin{tabular}{|c|p{3cm}|p{4cm}|p{8cm}|}
    \hline
      \textbf{Рік} & \textbf{Автор(и)} & \textbf{Підхід} & \textbf{Особливості} \\    
    \hline
      2023 & Ostadhossein \& Hoseinzadeh & МСР, МСО & Одновимірні лінійні і нелінійні задачі \\
    \hline
      2024 & B. Gangadhara et al. & МСЕ & Аналіз критеріїв стабільності \\
    \hline
      2023 & Santiago et al. & WRMG, МСР & Паралельний алгоритм, двовимірна академічна задача \\
    \hline
      2025 & Munshid et al. & МСЕ, BFO & Порівняння методів \\
    \hline
      2022 & Tian et al. & МСЕ, PBD & Інтеграція в симулятор iMSTK \\
    \hline
      2025 & Wu et al. & МСР, TLM & Розширена модель взаємодії хвиля - м'які тканини \\
    \hline
      2025 & Trefilov et al. & МСЕ, Аналітичний & Розширена модель для лазерної абляції \\
    \hline
  \end{tabular}
\end{center}

\begin{center}
  \textbf{Формулювання цілі статті}
\end{center}

Метою цієї праці є перевірка ефективності використання МПГЕ для розв'язування задач теплопровідності у людському тілі.
Отримані результати будуть використані для отримання рекомендацій щодо удосконалення методів моделювання теплових 
процесів на основі рівняння Пеннеса. Для проведення порівняльного аналізу було проаналізовано наукові статті, які 
описують різні методи розв'язування задач теплопровідності у людському тілі.

Для отримання результатів використовується МПГЕ описаний у \cite{grytsko-cobem} із різними товщинами приграничної зони.
Результати порівнюються на академічній задачі із відомим аналітичним розв'язком для кращого розуміння точності методів.
  
\begin{center}
  \textbf{Виклад основного матеріалу}
\end{center}

\begin{center}
  \textit{Постановка задачі}
\end{center}

Рівняння Пеннеса було сформульовано у 1948 році у статті \cite{thermal-process-tissue-pennes}. У честь п'ятдесятиріччя
оригінальної статті її параметри були уточнені і детальніше викладені \cite{thermal-process-tissue-pennes-explained}. На
даний момент у багатьох дослідженнях для опису моделі процесів всередині людського тіла використовують це рівняння.

У цій статті кома перед індексом означає похідну по змінній, яка відповідає цьому індексу, тобто ${f_{,j} =
\partial}{f/\partial}x_{j}$, ${{f_{,ij} = \partial^{2}}{f/\partial}x_{i}\partial x_{j}}{}$, при цьому використовується
німе сумування за індексами, які повторюються.

Задачу теплообміну у людських тканинах у заданій області $\Omega$ можна сформулювати у вигляді:

\begin{equation}
  \label{eqn:theory_thermo_pennes_eqn}
  \nabla \cdot \lambda_\theta \nabla \theta(x) + \omega \rho_b C_b (\theta_c - \theta(x)) + q_m  = 0, x \in \Omega
\end{equation}

\noindent де $\Gamma = \partial \Omega$ -- межа області $\Omega$, $\lambda_\theta$ -- коефіцієнт теплопровідності,
$\theta$ -- шукана температура в області $\Omega$,
$\rho_b$ -- густина крові, $\omega$ -- перфузія крові через орган, $C_b$ -- питома теплоємність крові,
$\theta_c$ -- внутрішня температура тіла, $q_m$ -- метаболічне тепло у тканині. На межі задамо умови:

\begin{equation}
  \label{eqn:theory_thermo_cond_dirichlet}
  \theta(x) = \theta_{\Gamma_D}(x), x \in \Gamma_D
\end{equation}

\begin{equation}
  \label{eqn:theory_thermo_cond_neumann}
  q(x) = q_{\Gamma_N}(x), x \in \Gamma_N
\end{equation}

\begin{equation}
  \label{eqn:theory_thermo_cond_robin}
  -\lambda_\theta q(x) = h(\theta(x) - \theta_e), x \in \Gamma_R,
\end{equation}

\noindent де $q(x) = \frac{\partial \theta(x)}{\partial n}$ -- тепловий потік, $n_i$ -- компоненти вектора зовнішньої
нормалі до межі $\Gamma$, $\theta_e$ -- температура середовища, $h$ -- коефіцієнт тепловіддачі. Крайові умови
(\ref{eqn:theory_thermo_cond_dirichlet}), (\ref{eqn:theory_thermo_cond_neumann}) і (\ref{eqn:theory_thermo_cond_robin})
-- це умови Діріхле, Неймана і Робіна відповідно. Умова Робіна описує тепловий обмін тіла із навколишнім середовищем.

Для розв'язування задач типу (\ref{eqn:theory_thermo_pennes_eqn}) (\ref{eqn:theory_thermo_cond_dirichlet})
(\ref{eqn:theory_thermo_cond_neumann}) (\ref{eqn:theory_thermo_cond_robin}) використаємо метод приграничних елементів
(МПГЕ)\cite{grytsko-cobem}, який належить до сімейства методів граничних елементів (МГЕ)\cite{brebbia-bem}. МГЕ мають
кілька переваг у порівнянні із методом скінчених елементів (МСЕ)\cite{bahvalov-et-al}, зокрема:
\begin{itemize}
  \item Розбиття границі області на дискретні елементи, а не цілої область.
  \item Зменшує розмірність задачі і відповідної кількості лінійних алгебраїчних рівнянь.
  \item Висока точність розв'язку на границі області.
\end{itemize}
Такі переваги роблять МГЕ привабливими для використання у симуляторах.

МПГЕ дискретизує не саму границю, а приграничну зону. Це дозволяє ввести додатковий ступінь вільності, а саме ширину
цієї зони. Він дозволяє покращувати точність стаціонарних задач. Також інтеграли, які обчислюються мають сингулярність
на межі, а не всередині області інтегрування. Це дозволяє використовувати звичайні формули інтегрування для обчислення
розв'язку. Симулятори можуть базуватися на МПГЕ.

\begin{center}
  \textit{Інтегральне подання}
\end{center}

Візьмемо фундаментальні розв'язки рівняння Пеннеса (\ref{eqn:theory_thermo_pennes_eqn}) із
\cite{pennes-equation-bem-fundamental-solutions} і позначимо їх $\hat{\theta}$ і $\hat{q}$: 

\begin{equation*}
  \hat{\theta}(x, y) = \frac{K_0(r(x, y) \sqrt{\omega \rho_b C_b})}{2\pi}
\end{equation*}

\begin{equation*}
  \hat{q}(x, y) = \frac{\partial \hat{\theta}(x, y)}{\partial n},
\end{equation*}

\noindent де $K_0(x)$ -- модифікована функція Бесселя другого роду нульового порядку, $r(x, y)$ -- відстань між точками
x і y, а $n$ -- нормаль.

Розглянемо таку область $B \subset R^n$, що $\Omega \subset B$, $\partial\Omega \cap \partial B = \varnothing$.
  Позначимо приграничну область $G = B \setminus \Omega$. Інтегральне подання розв'язку для МПГЕ запишемо у вигляді:

\begin{equation}
  \label{eqn:thermo_cobem_pennes_theta}
  \theta(\xi) = \int \limits_G g(x)\hat{\theta}(x,\xi)dG(x)
\end{equation}

\begin{equation}
  \label{eqn:thermo_cobem_pennes_q}
  q(\xi) = \int \limits_G g(x)\hat{q}(x,\xi)dG(x)
\end{equation}

Область $G$ розіб'ємо на приграничні елементи $G_j$, $j=\overline{1,n}$. Кожен із цих елементів містить кусок границі
області $\Omega$, $G_j \cap \Gamma \neq \varnothing$. Визначимо область $G$, яка містить область $\Omega$, причому
$\partial G \cap \partial \Omega = \varnothing$.

\begin{center}
  \textit{Дискредитація}
\end{center}

Дискретизуємо інтегральне подання \ref{eqn:thermo_cobem_pennes_theta}:

\begin{equation}
  \label{eqn:thermo_cobem_pennes_theta_discrete}
  \theta(\xi) = \sum_j^n g_j  \int \limits_{G_j} \hat{\theta}(x,\xi)dG_j(x),
\end{equation}

\noindent де $g_j$ -- стале значення функції $g(x)$ на усьому пригранич­ному елементі $G_j$. Відповідно у точці
колокації $x_i$ отримаємо:

\begin{equation*}
  \theta(x_i) = \sum_j^n g_j \int \limits_{G_j} \hat{\theta}(x,x_i)dG_j(x)
\end{equation*}

\noindent Таке рівняння можна переписати у вигляді $\theta_i = \theta_{\Gamma_D}(x_i) = H_{ij} g_j$, де: 

\begin{equation*}
  H_{ij} = \sum_j^n \int \limits_{G_j} \hat{\theta}(x,x_i)dG_j(x)
\end{equation*}

\noindent Відповідно, $q_i = q_{\Gamma_N}(x_i) = P_{ij} g_j$ -- це дискретне представлення для інтегрального подання
(\ref{eqn:thermo_cobem_pennes_q}) і умови (\ref{eqn:theory_thermo_cond_neumann}), де:

\begin{equation*}
  P_{ij} = \int \limits_{G_j} \hat{q}(x,x_i)dG_j(x)
\end{equation*}

Система лінійних алгебраїчних рівнянь (СЛАР) запишемо у вигляді:

\begin{equation*}
  \begin{bmatrix}
    H_{11} & ... & H_{1k} & ... & H_{1l} & ...  & H_{1n} \\
    ... \\
    H_{k1} & ... & H_{kk} & ... & H_{kl} & ...  & H_{kn} \\
    ... \\
    P_{l1} & ... & P_{lk} & ... & P_{ll} & ...  & P_{ln} \\
    ... \\
    P_{n1} & ... & P_{nk} & ... & P_{nl} & ...  & P_{nn}
  \end{bmatrix}
  \begin{bmatrix}
    g_1 \\
    ...   \\
    g_k \\
    ... \\
    g_l \\
    ... \\
    g_n
  \end{bmatrix}
  =
  \begin{bmatrix}
    \theta_1 \\
    ...   \\
    \theta_k \\
    ... \\
    q_l \\
    ... \\
    q_n
  \end{bmatrix}
\end{equation*}

\noindent Розв'язок СЛАР підставимо у рівняння (\ref{eqn:thermo_cobem_pennes_theta_discrete}). Завдяки цьому ми зможемо
обчислити температуру у довільній точці області $\Omega$.

\begin{center}
  \textbf{Порівняння методів}
\end{center}

Для порівняння методів візьмемо область $\Omega$ у вигляді квадратної пластини, довжина сторони 1 сантиметр.
Розв'яжемо задачу (\ref{eqn:theory_thermo_pennes_eqn}) (\ref{eqn:theory_thermo_cond_dirichlet}). Візьмемо наступні
значення коефіцієнтів: $\lambda_{\theta} = 1 \frac{\text{Вт}}{\text{м}\cdot\text{К}}$, $\omega = 1 \frac{1}{c}$, $\rho_b
= 1 \frac{\text{кг}}{\text{м}^3}$, $C_b = 2 \frac{\text{Дж}}{\text{кг}\cdot\text{К}}$, $\theta_c = {0}^\circ C$ і $q_m =
0 \frac{\text{Вт}}{\text{м}^3}$. На границі області $\Omega$ задамо значення відомого аналітичного розв'язку 
$\theta(x) = sinh(x_2)$.

Порівняємо результати отримані МСЕ і МПГЕ. У першому чисельному експерименті було використано 56 скінчених елементів і
40 приграничних елементів, ширина приграничної смуги 0,4. Максимальна похибка у МСЕ приблизно рівна 0,09, а в МПГЕ 0,06. 

\begin{figure}[ht!]
  \centering
  \begin{subfigure}{.5\textwidth}
    \centering
    \includegraphics[width=.8\linewidth]{fem_56_error.png}
    \caption{МСЕ (56 елементів)}
    \label{fig:experiment_1_fem}
  \end{subfigure}%
  \begin{subfigure}{.5\textwidth}
    \centering
    \includegraphics[width=.8\linewidth]{cobem_40_0_4_error.png}
    \caption{МПГЕ (40 елементів)}
    \label{fig:experiment_1_cobem}
  \end{subfigure}
  \label{fig:experiment_1}
\end{figure}

У другому експерименті було 226 скінчених і 120 приграничних елементів. У МСЕ максимальна похибка приблизно рівна 0,09, а в МПГЕ -- 0,00025. В обох
чисельних експериментах ширина приграничної смуги 0,4.

\begin{figure}[ht!]
  \centering
  \begin{subfigure}{.5\textwidth}
    \centering
    \includegraphics[width=.8\linewidth]{fem_226_error.png}
    \caption{МСЕ (226 елементів)}
    \label{fig:experiment_2_fem}
  \end{subfigure}%
  \begin{subfigure}{.5\textwidth}
    \centering
    \includegraphics[width=.8\linewidth]{cobem_120_0_4_error.png}
    \caption{МПГЕ (120 елементів)}
    \label{fig:experiment_2_cobem}
  \end{subfigure}
  \label{fig:experiment_2}
\end{figure}

\begin{center}
  \textbf{Подальші дослідження}
\end{center}

У статті наведені приклади розв'язування стаціонарної задачі для рівняння Пеннеса. Всі коефіцієнти цього рівняння були 
сталими і не залежали від температури чи координат. Тому одним з кроків у подальших дослідженнях є розгляд неоднорідної
поведінки фізичного середовища у певній області, яка залежить від координат. Це буде відповідати дослідженням 
пов'язаними із описом фізичних властивостей тканин людського тіла в області злоякісних включень. Також вартим уваги є 
задача у тривимірному варіанті. Наступним етапом є розгляд задач із кількома злоякісними включеннями, які описуються 
неоднорідностями коефіцієнтів рівняння Пеннеса, що відповідають за опис фізичних властивостей середовища.

\begin{center}
  \textbf{Висновки}
\end{center}
  
Існує значна кількість методів і підходів для моделювання процесів теплового обміну в організмі людині. Це спостереження
узгоджується із аналізом предметної області проведеним у цій статті. Проте розвиток медичної індустрії потребує
постійного удосконалення засобів моделювання фізичних явищ в тілі людини. Зокрема пришвидшення процесу моделювання для
ефективнішої симуляції теплового обміну в тілі людини.

Проаналізовано наукові праці, які розв'язують рівняння Пеннеса. На основі цього аналізу створено систематизацію методів,
які використовуються на практиці. Серед цим методів традиційно використовується МСЕ, який є класичним чисельним методом
для багатьох задач математичного моделювання фізичних процесів. Його вибрано для порівняльного аналізу із МПГЕ.

Досліджено фундаментальні розв'язки для рівняння Пеннеса. Застосовано МПГЕ для розв'язування рівняння Пеннеса. Зроблено
порівняльний аналіз між МСЕ і МПГЕ для задачі із відомим аналітичним розв'язком. Показано ефективність запропонованого
підходу. Зокрема використання ширини приграничної області для підвищення точності обчислень. Також, потрібну меншу
кількість приграничних елементів для отримання результатів із високою точністю. Запропоновано подальші кроки в
удосконаленні математичних моделей.   
